% GENERATED FILE. Edit canonical lessons, then run npm run generate:books.
% canonical-source-hash: fnv1a64:968e5202a1e8183f
% canonical-lessons: SA-C224-kendram, SA-C224-attalika, SA-C224-andhakarah, SA-C224-jyotih, SA-C224-bhumih

\chapter{Centre, Tall building, Darkness, Radiance, Ground}
\label{ch:sa-centre-tall-building-darkness-radiance-ground}
\hlchaptermodality{224}

\begin{chapteropening}
\textbf{By the end of this chapter:} \emph{I can say a centre, a tall building, darkness, radiance and the ground.}

Complete the last lesson of chapter 224: I can say a centre, a tall building, darkness, radiance and the ground.
\end{chapteropening}

\section[\texorpdfstring{\sk{केन्द्रम्} (kendram)}{केन्द्रम् (kendram)}]{\sk{केन्द्रम्} (kendram) — a centre}

\label{lesson:SA-C224-kendram}



\begin{quote}
Before the new one: say the Sanskrit for a hermitage, then the Sanskrit for a monastery.
\end{quote}

\subsection*{You'll want to know: \sk{केन्द्रम्}}
\textbf{\sk{केन्द्रम्}} — \emph{kendram} — \textquotedblleft{}a centre\textquotedblright{}.

\begin{grammarlens}[title={What you've built}]
One more word for this chapter.
\end{grammarlens}

\subsection*{Your turn}
\begin{itemize}
\raggedright
  \item \emph{Say it:} \emph{kendram}
  \item \emph{Say it:} \emph{kendram}, once more
  \item \emph{Say it:} say \emph{kendram}
  \item \emph{Recall:} say the Sanskrit for a hermitage, then the Sanskrit for a monastery, then say \emph{kendram} again
\end{itemize}

\begin{tcolorbox}[breakable,colback=teal!4,colframe=teal!35!black,title={Before you move on}]
What does \sk{केन्द्रम्} mean? (\textquotedblleft{}A centre\textquotedblright{}.) Say it once more.
\end{tcolorbox}
\section[\texorpdfstring{\sk{अट्टालिका} (aṭṭālikā)}{अट्टालिका (aṭṭālikā)}]{\sk{अट्टालिका} (aṭṭālikā) — a tall building}

\label{lesson:SA-C224-attalika}



\begin{quote}
Before the new one: say the Sanskrit for a monastery, then the Sanskrit for a centre.
\end{quote}

\subsection*{You'll want to know: \sk{अट्टालिका}}
\textbf{\sk{अट्टालिका}} — \emph{aṭṭālikā} — \textquotedblleft{}a tall building\textquotedblright{}.

\begin{grammarlens}[title={What you've built}]
One more word for this chapter.
\end{grammarlens}

\subsection*{Your turn}
\begin{itemize}
\raggedright
  \item \emph{Say it:} \emph{aṭṭālikā}
  \item \emph{Say it:} \emph{aṭṭālikā}, once more
  \item \emph{Say it:} say \emph{aṭṭālikā}
  \item \emph{Recall:} say the Sanskrit for a monastery, then the Sanskrit for a centre, then say \emph{aṭṭālikā} again
\end{itemize}

\begin{tcolorbox}[breakable,colback=teal!4,colframe=teal!35!black,title={Before you move on}]
What does \sk{अट्टालिका} mean? (\textquotedblleft{}A tall building\textquotedblright{}.) Say it once more.
\end{tcolorbox}
\section[\texorpdfstring{\sk{अन्धकारः} (andhakāraḥ)}{अन्धकारः (andhakāraḥ)}]{\sk{अन्धकारः} (andhakāraḥ) — darkness}

\label{lesson:SA-C224-andhakarah}



\begin{quote}
Before the new one: say the Sanskrit for a centre, then the Sanskrit for a tall building.
\end{quote}

\subsection*{You'll want to know: \sk{अन्धकारः}}
\textbf{\sk{अन्धकारः}} — \emph{andhakāraḥ} — \textquotedblleft{}darkness\textquotedblright{}.

\begin{grammarlens}[title={What you've built}]
One more word for this chapter.
\end{grammarlens}

\subsection*{Your turn}
\begin{itemize}
\raggedright
  \item \emph{Say it:} \emph{andhakāraḥ}
  \item \emph{Say it:} \emph{andhakāraḥ}, once more
  \item \emph{Say it:} say \emph{andhakāraḥ}
  \item \emph{Recall:} say the Sanskrit for a centre, then the Sanskrit for a tall building, then say \emph{andhakāraḥ} again
\end{itemize}

\begin{tcolorbox}[breakable,colback=teal!4,colframe=teal!35!black,title={Before you move on}]
What does \sk{अन्धकारः} mean? (\textquotedblleft{}Darkness\textquotedblright{}.) Say it once more.
\end{tcolorbox}
\section[\texorpdfstring{\sk{ज्योतिः} (jyotiḥ)}{ज्योतिः (jyotiḥ)}]{\sk{ज्योतिः} (jyotiḥ) — radiance, a flame}

\label{lesson:SA-C224-jyotih}



\begin{quote}
Before the new one: say the Sanskrit for a tall building, then the Sanskrit for darkness.
\end{quote}

\subsection*{You'll want to know: \sk{ज्योतिः}}
\textbf{\sk{ज्योतिः}} — \emph{jyotiḥ} — \textquotedblleft{}radiance, a flame\textquotedblright{}.

\begin{grammarlens}[title={What you've built}]
One more word for this chapter.
\end{grammarlens}

\subsection*{Your turn}
\begin{itemize}
\raggedright
  \item \emph{Say it:} \emph{jyotiḥ}
  \item \emph{Say it:} \emph{jyotiḥ}, once more
  \item \emph{Say it:} say \emph{jyotiḥ}
  \item \emph{Recall:} say the Sanskrit for a tall building, then the Sanskrit for darkness, then say \emph{jyotiḥ} again
\end{itemize}

\begin{tcolorbox}[breakable,colback=teal!4,colframe=teal!35!black,title={Before you move on}]
What does \sk{ज्योतिः} mean? (\textquotedblleft{}Radiance, a flame\textquotedblright{}.) Say it once more.
\end{tcolorbox}
\section[\texorpdfstring{\sk{भूमिः} (bhūmiḥ)}{भूमिः (bhūmiḥ)}]{\sk{भूमिः} (bhūmiḥ) — the ground, land}

\label{lesson:SA-C224-bhumih}



\begin{quote}
Before the new one: say the Sanskrit for darkness, then the Sanskrit for radiance.
\end{quote}

\subsection*{You'll want to know: \sk{भूमिः}}
\textbf{\sk{भूमिः}} — \emph{bhūmiḥ} — \textquotedblleft{}the ground, land\textquotedblright{}.

\begin{grammarlens}[title={What you've built}]
One more word for this chapter.
\end{grammarlens}

\subsection*{Your turn}
\begin{itemize}
\raggedright
  \item \emph{Say it:} \emph{bhūmiḥ}
  \item \emph{Say it:} \emph{bhūmiḥ}, once more
  \item \emph{Say it:} say \emph{bhūmiḥ}
  \item \emph{Recall:} say the Sanskrit for darkness, then the Sanskrit for radiance, then say \emph{bhūmiḥ} again
\end{itemize}

\begin{tcolorbox}[breakable,colback=teal!4,colframe=teal!35!black,title={Before you move on}]
What does \sk{भूमिः} mean? (\textquotedblleft{}The ground, land\textquotedblright{}.) Say it once more.
\end{tcolorbox}
